\chapter{Random Walks}

A drunk man leaves a bar at midnight. He can barely stand, and certainly
cannot navigate. Each step he takes is a coin flip: east or west, with no
memory of where he has been and no plan for where he is going. Will he ever
stumble home? Or will he wander off into the night, never to return?

This is not an idle riddle. It is one of the most fertile questions in all of
probability, and the path from it leads -- astonishingly directly -- to the
physics of diffusing molecules, the price of a stock option, the spread of a
gene through a population, the design of search algorithms, and the
mathematics of Brownian motion. The humble \emph{random walk} is the bridge
between the discrete world of coin flips, which we have studied at length, and
the continuous world of motion and noise.

In this chapter we build the random walk from the ground up. We start with the
simplest possible version -- steps of $\pm 1$ on a line -- and we ask the
natural questions a curious person asks: where do you end up on average, how
far do you spread, do you come back to the start, and what does a typical path
actually look like? The answers are sometimes intuitive, sometimes deeply
surprising, and always illuminating.

\chapterline

\section{The Simple Random Walk}
\label{ch26:srw}

\begin{intuitionbox}
Forget the formalism for a moment. Picture a token sitting on the integer $0$
of a number line. A fair coin is flipped. Heads, the token moves one step
right; tails, one step left. Flip again. And again. The position of the token
after many flips is a random walk. Everything we do in this chapter is a
question about that wandering token.
\end{intuitionbox}

Let $X_1, X_2, X_3, \dots$ be independent steps, each equal to $+1$ or $-1$
with probability $1/2$:
\[
  \P(X_i = +1) = \P(X_i = -1) = \tfrac{1}{2}, \qquad i = 1, 2, 3, \dots
\]
These are independent and identically distributed (i.i.d.) -- the coin has no
memory. The \emph{position} of the walk after $n$ steps is the running sum
\[
  S_n = X_1 + X_2 + \cdots + X_n, \qquad S_0 = 0.
\]

\begin{definition}[Simple symmetric random walk]
\label{ch26:def-srw}
The sequence $S_0, S_1, S_2, \dots$ defined by $S_0 = 0$ and
$S_n = \sum_{i=1}^n X_i$, where the $X_i$ are i.i.d.\ with
$\P(X_i = \pm 1) = 1/2$, is the \emph{simple symmetric random walk} on $\Z$.
If instead $\P(X_i = +1) = p$ and $\P(X_i = -1) = q = 1-p$ with $p \neq 1/2$,
we call it a \emph{biased} (or asymmetric) random walk.
\end{definition}

The word ``simple'' means the steps have size one; ``symmetric'' means the coin
is fair. We will spend most of our energy on the symmetric case, because that
is where the deepest surprises live, but the biased walk is exactly what we
need for gambling problems, so we keep it in our pocket.

\begin{happeningbox}
A random walk is a \emph{cumulative sum of i.i.d.\ random variables}. That one
sentence connects this entire chapter to everything you already know. The
Law of Large Numbers and the Central Limit Theorem are theorems about sums
$S_n$ -- so they are secretly theorems about random walks. When we ask ``where
is the walk after $n$ steps?'' we are asking about the distribution of a sum,
and the CLT will hand us a bell curve. The new questions -- recurrence,
first-passage times, the reflection principle -- are the ones that go
\emph{beyond} a single snapshot $S_n$ and ask about the whole \emph{path}.
\end{happeningbox}

\subsection{A first look at a path}

Before any formulas, let us just \emph{watch}. Figure~\ref{ch26:fig-paths}
shows several independent sample paths of the symmetric walk over $120$ steps.
Stare at them. Notice three things. First, the paths sprawl: by step $120$ they
are scattered over a wide band, not clustered near zero. Second, they are
jagged and unpredictable -- no path ``trends.'' Third, and most importantly,
the band of paths widens like the square root of time, not linearly. We will
prove that shape precisely in the next section, but your eye can already see
it.

\begin{figure}[H]
\centering
\begin{tikzpicture}
\begin{axis}[
  width=13cm, height=7cm,
  xlabel={step $n$}, ylabel={position $S_n$},
  xmin=0, xmax=120, ymin=-26, ymax=26,
  grid=both, grid style={gray!18},
  title={Five sample paths of the simple symmetric random walk},
  legend style={font=\scriptsize, at={(1.02,1)}, anchor=north west}
]
\addplot[sourceblue, thick] coordinates {
(0,0)(1,1)(2,0)(3,1)(4,2)(5,1)(6,2)(7,3)(8,2)(9,3)(10,4)(11,3)(12,4)(13,5)(14,4)(15,5)(16,6)(17,5)(18,6)(19,7)(20,6)(21,7)(22,8)(23,7)(24,6)(25,7)(26,8)(27,9)(28,8)(29,9)(30,10)(31,9)(32,10)(33,11)(34,10)(35,9)(36,10)(37,11)(38,12)(39,11)(40,12)(41,13)(42,12)(43,13)(44,12)(45,13)(46,14)(47,13)(48,12)(49,13)(50,14)(51,15)(52,14)(53,15)(54,16)(55,15)(56,14)(57,15)(58,16)(59,15)(60,16)(61,17)(62,16)(63,15)(64,16)(65,17)(66,16)(67,17)(68,18)(69,17)(70,16)(71,17)(72,16)(73,17)(74,18)(75,17)(76,16)(77,15)(78,16)(79,17)(80,18)(81,17)(82,18)(83,19)(84,18)(85,17)(86,18)(87,17)(88,18)(89,19)(90,18)(91,17)(92,18)(93,19)(94,20)(95,19)(96,18)(97,19)(98,18)(99,19)(100,20)(101,19)(102,18)(103,19)(104,18)(105,17)(106,18)(107,19)(108,18)(109,17)(110,18)(111,19)(112,18)(113,17)(114,18)(115,19)(116,18)(117,17)(118,18)(119,17)(120,18)
};
\addplot[densityred, thick] coordinates {
(0,0)(1,-1)(2,-2)(3,-1)(4,-2)(5,-3)(6,-2)(7,-3)(8,-4)(9,-3)(10,-2)(11,-3)(12,-2)(13,-1)(14,-2)(15,-1)(16,-2)(17,-3)(18,-2)(19,-1)(20,-2)(21,-1)(22,0)(23,-1)(24,-2)(25,-1)(26,0)(27,-1)(28,-2)(29,-3)(30,-2)(31,-3)(32,-4)(33,-5)(34,-4)(35,-5)(36,-6)(37,-5)(38,-6)(39,-7)(40,-6)(41,-7)(42,-8)(43,-7)(44,-8)(45,-9)(46,-8)(47,-9)(48,-10)(49,-9)(50,-10)(51,-11)(52,-10)(53,-11)(54,-10)(55,-11)(56,-12)(57,-11)(58,-12)(59,-13)(60,-12)(61,-11)(62,-12)(63,-13)(64,-12)(65,-11)(66,-12)(67,-13)(68,-14)(69,-13)(70,-14)(71,-15)(72,-14)(73,-15)(74,-16)(75,-15)(76,-14)(77,-15)(78,-14)(79,-13)(80,-14)(81,-15)(82,-14)(83,-13)(84,-14)(85,-13)(86,-12)(87,-13)(88,-12)(89,-13)(90,-14)(91,-13)(92,-12)(93,-13)(94,-12)(95,-13)(96,-14)(97,-13)(98,-12)(99,-11)(100,-12)(101,-13)(102,-12)(103,-13)(104,-14)(105,-13)(106,-14)(107,-15)(108,-14)(109,-15)(110,-16)(111,-15)(112,-16)(113,-17)(114,-16)(115,-17)(116,-16)(117,-15)(118,-16)(119,-15)(120,-14)
};
\addplot[targetgreen, thick] coordinates {
(0,0)(1,1)(2,2)(3,1)(4,0)(5,1)(6,0)(7,-1)(8,0)(9,1)(10,0)(11,1)(12,2)(13,1)(14,0)(15,-1)(16,0)(17,1)(18,0)(19,-1)(20,0)(21,1)(22,2)(23,3)(24,2)(25,1)(26,2)(27,3)(28,2)(29,1)(30,2)(31,3)(32,4)(33,3)(34,2)(35,3)(36,4)(37,5)(38,4)(39,3)(40,4)(41,5)(42,4)(43,5)(44,6)(45,5)(46,6)(47,7)(48,6)(49,5)(50,4)(51,5)(52,6)(53,5)(54,4)(55,3)(56,4)(57,5)(58,4)(59,3)(60,2)(61,3)(62,2)(63,1)(64,2)(65,3)(66,2)(67,1)(68,0)(69,1)(70,2)(71,1)(72,2)(73,3)(74,4)(75,3)(76,2)(77,3)(78,4)(79,5)(80,4)(81,3)(82,2)(83,1)(84,2)(85,3)(86,4)(87,3)(88,2)(89,1)(90,0)(91,-1)(92,0)(93,1)(94,2)(95,1)(96,0)(97,1)(98,2)(99,3)(100,2)(101,3)(102,4)(103,3)(104,2)(105,1)(106,2)(107,3)(108,2)(109,1)(110,0)(111,1)(112,2)(113,3)(114,2)(115,1)(116,0)(117,1)(118,0)(119,-1)(120,0)
};
\addplot[quotientorange, thick] coordinates {
(0,0)(1,-1)(2,0)(3,1)(4,0)(5,-1)(6,-2)(7,-1)(8,0)(9,1)(10,2)(11,3)(12,2)(13,3)(14,4)(15,5)(16,4)(17,3)(18,4)(19,5)(20,6)(21,5)(22,6)(23,7)(24,8)(25,7)(26,8)(27,9)(28,10)(29,9)(30,10)(31,11)(32,10)(33,11)(34,12)(35,11)(36,12)(37,13)(38,12)(39,11)(40,12)(41,13)(42,14)(43,13)(44,14)(45,15)(46,14)(47,13)(48,14)(49,15)(50,16)(51,15)(52,16)(53,17)(54,16)(55,17)(56,18)(57,19)(58,18)(59,17)(60,18)(61,19)(62,20)(63,19)(64,20)(65,21)(66,20)(67,21)(68,20)(69,21)(70,22)(71,21)(72,20)(73,21)(74,22)(75,23)(76,22)(77,21)(78,22)(79,21)(80,22)(81,23)(82,22)(83,21)(84,22)(85,23)(86,22)(87,21)(88,22)(89,23)(90,22)(91,23)(92,22)(93,21)(94,22)(95,21)(96,20)(97,21)(98,22)(99,21)(100,22)(101,21)(102,20)(103,21)(104,20)(105,21)(106,22)(107,21)(108,20)(109,21)(110,22)(111,21)(112,20)(113,21)(114,22)(115,21)(116,22)(117,21)(118,20)(119,21)(120,22)
};
\addplot[structurepurple, thick] coordinates {
(0,0)(1,1)(2,0)(3,-1)(4,-2)(5,-1)(6,0)(7,-1)(8,-2)(9,-1)(10,-2)(11,-1)(12,0)(13,1)(14,0)(15,1)(16,2)(17,1)(18,0)(19,-1)(20,-2)(21,-1)(22,-2)(23,-3)(24,-4)(25,-3)(26,-2)(27,-3)(28,-4)(29,-5)(30,-4)(31,-3)(32,-2)(33,-3)(34,-4)(35,-5)(36,-6)(37,-5)(38,-6)(39,-7)(40,-6)(41,-5)(42,-4)(43,-5)(44,-6)(45,-7)(46,-8)(47,-7)(48,-8)(49,-9)(50,-8)(51,-7)(52,-8)(53,-9)(54,-8)(55,-7)(56,-6)(57,-7)(58,-8)(59,-9)(60,-8)(61,-9)(62,-10)(63,-9)(64,-8)(65,-7)(66,-8)(67,-9)(68,-10)(69,-9)(70,-8)(71,-9)(72,-10)(73,-11)(74,-10)(75,-9)(76,-10)(77,-11)(78,-12)(79,-11)(80,-10)(81,-9)(82,-10)(83,-11)(84,-10)(85,-9)(86,-8)(87,-9)(88,-10)(89,-11)(90,-10)(91,-9)(92,-8)(93,-9)(94,-8)(95,-7)(96,-8)(97,-9)(98,-10)(99,-9)(100,-8)(101,-9)(102,-10)(103,-11)(104,-10)(105,-11)(106,-12)(107,-11)(108,-10)(109,-9)(110,-10)(111,-11)(112,-10)(113,-11)(114,-12)(115,-11)(116,-10)(117,-9)(118,-10)(119,-11)(120,-10)
};
\addplot[black, dashed, thick, samples=80, domain=0:120] {sqrt(x)};
\addplot[black, dashed, thick, samples=80, domain=0:120] {-sqrt(x)};
\legend{path 1, path 2, path 3, path 4, path 5}
\end{axis}
\end{tikzpicture}
\caption{Five independent sample paths of the simple symmetric random walk.
The dashed curves are $\pm\sqrt{n}$: the typical spread of the walk grows like
the \emph{square root} of the number of steps, not linearly. Most paths stay
within a band of this width, though any single path can wander outside it.}
\label{ch26:fig-paths}
\end{figure}

\section{Expected Position and Variance}
\label{ch26:moments}

Where is the walk, on average, after $n$ steps? The answer feels almost
like a trick.

\begin{theorem}[Mean and variance of the walk]
\label{ch26:thm-meanvar}
For the simple symmetric random walk,
\[
  \E[S_n] = 0, \qquad \Var(S_n) = n, \qquad \SD(S_n) = \sqrt{n}.
\]
For the biased walk with $\P(X_i = +1) = p$, $\P(X_i = -1) = q = 1-p$, writing
the drift $\mu = p - q$ and $\sigma^2 = 1 - \mu^2$,
\[
  \E[S_n] = n\mu = n(p-q), \qquad \Var(S_n) = n\sigma^2 = 4npq.
\]
\end{theorem}

\begin{whybox}
Both lines follow from a single idea: expectation is linear, and for
independent variables variances add.

\textbf{The mean.} Each step has $\E[X_i] = (+1)p + (-1)q = p - q = \mu$. By
linearity of expectation -- which needs no independence at all --
\[
  \E[S_n] = \E\!\Big[\sum_{i=1}^n X_i\Big] = \sum_{i=1}^n \E[X_i] = n\mu.
\]
For the fair coin $p = q = 1/2$, so $\mu = 0$ and $\E[S_n] = 0$.

\textbf{The variance.} First, a single step. Since $X_i \in \{-1,+1\}$ we
have $X_i^2 = 1$ always, so $\E[X_i^2] = 1$ and
\[
  \Var(X_i) = \E[X_i^2] - (\E[X_i])^2 = 1 - \mu^2 = \sigma^2.
\]
For the fair coin $\mu = 0$ so $\Var(X_i) = 1$. Now because the steps are
\emph{independent}, the variance of the sum is the sum of the variances:
\[
  \Var(S_n) = \sum_{i=1}^n \Var(X_i) = n\sigma^2.
\]
For the fair coin this is just $n$. Taking square roots gives the standard
deviation $\sqrt{n}\,\sigma$, and $\sqrt{n}$ in the fair case.
\end{whybox}

That $\E[S_n] = 0$ is the trick the symmetric walk plays on your intuition. On
average, the drunk goes nowhere. But ``on average zero'' does not mean ``stays
near zero.'' It means the rightward and leftward excursions cancel \emph{in
expectation}, the way the average of $+1000$ and $-1000$ is zero. The right
measure of how far the walk roams is not the mean but the standard deviation,
and that grows without bound, like $\sqrt{n}$.

\begin{focusbox}
\textbf{The $\sqrt{n}$ law.} After $n$ steps the walk is typically a distance
of order $\sqrt{n}$ from the origin. After $100$ steps, about $10$ away; after
$10{,}000$ steps, about $100$ away. To get twice as far you must walk four
times as long. This sublinear, square-root spreading is the signature of
diffusion, and it reappears verbatim in the physics of heat, ink in water, and
pollen grains under a microscope.
\end{focusbox}

\begin{example}[How far is ``far''?]
\label{ch26:ex-howfar}
A particle takes $n = 10{,}000$ symmetric unit steps. What is the probability
it ends up more than $300$ units from the origin?

Here $\Var(S_n) = 10{,}000$, so $\SD(S_n) = 100$. A displacement of $300$ is
exactly $3$ standard deviations. By the Central Limit Theorem (the sum of
$10{,}000$ i.i.d.\ bounded steps is very nearly normal), $S_n \approx
\Normal(0, 100^2)$, and
\[
  \P(\abs{S_n} > 300) \approx \P(\abs{Z} > 3) \approx 0.0027,
\]
about a quarter of one percent. Ending up $3000$ away (thirty standard
deviations) is, for all practical purposes, impossible -- even though the walk
\emph{could} in principle go straight right the entire time, with probability
$2^{-10000}$.
\end{example}

\begin{connectionbox}
The exact distribution of $S_n$ is a relabelled binomial. If $R$ is the number
of rightward steps, then $R \sim \Bin(n, p)$ and $S_n = R - (n - R) = 2R - n$.
So
\[
  \P(S_n = k) = \P\!\Big(R = \tfrac{n+k}{2}\Big)
  = \binom{n}{\frac{n+k}{2}} p^{\frac{n+k}{2}} q^{\frac{n-k}{2}},
\]
valid when $n + k$ is even (the walk's parity must match $n$: after an even
number of steps you sit on an even site). Everything about the position $S_n$
is binomial in disguise -- including its bell shape, which is the de
Moivre--Laplace theorem, the historical seed of the CLT. See the companion
\emph{Calculus From the Ground Up} for how that bell shape emerges from
Stirling's approximation to the binomial coefficient.
\end{connectionbox}

Figure~\ref{ch26:fig-pmf} shows this binomial pmf for $n = 40$ together with
the normal curve it converges to. The match is already excellent.

\begin{figure}[H]
\centering
\begin{tikzpicture}
\begin{axis}[
  width=13cm, height=6.2cm,
  xlabel={position $S_{40}$}, ylabel={$\P(S_{40}=k)$},
  xmin=-22, xmax=22, ymin=0, ymax=0.14,
  ybar, bar width=4pt,
  title={Distribution of position after $n=40$ steps (bars) vs.\ its normal approximation},
]
\addplot[draw=sourceblue, fill=histfill] coordinates {
(-20,0.0008)(-18,0.0021)(-16,0.0051)(-14,0.0109)(-12,0.0211)(-10,0.0366)(-8,0.0572)(-6,0.0807)(-4,0.1031)(-2,0.1194)(0,0.1254)(2,0.1194)(4,0.1031)(6,0.0807)(8,0.0572)(10,0.0366)(12,0.0211)(14,0.0109)(16,0.0051)(18,0.0021)(20,0.0008)
};
\addplot[densityred, very thick, samples=160, domain=-22:22]
  {2*1/(sqrt(2*pi)*sqrt(40))*exp(-x^2/(2*40))};
\end{axis}
\end{tikzpicture}
\caption{The exact distribution of $S_{40}$ (only even sites have positive
probability, since parity must match $n=40$) overlaid with the normal density
$\Normal(0,40)$. The bars are spaced two apart, so each carries the
probability of a width-$2$ interval; the curve is scaled to match. The CLT is
visible to the naked eye.}
\label{ch26:fig-pmf}
\end{figure}

\section{Recurrence and Transience: Does the Drunk Get Home?}
\label{ch26:recurrence}

Now the question that gives this chapter its soul. The walk wanders. Does it
ever come \emph{back}? More sharply: with probability $1$, will the walk return
to its starting point $0$?

\begin{definition}[Recurrence and transience]
\label{ch26:def-recurrence}
A random walk is \emph{recurrent} if it returns to its starting point with
probability $1$ (and hence, by starting afresh each time, returns infinitely
often). It is \emph{transient} if there is positive probability it never
returns -- in which case it visits the origin only finitely many times and
eventually drifts away forever.
\end{definition}

The answer in one dimension is one of the loveliest facts in probability.

\begin{theorem}[Recurrence of the 1D symmetric walk]
\label{ch26:thm-rec1d}
The simple symmetric random walk on $\Z$ is recurrent: it returns to the
origin with probability $1$. In fact it visits \emph{every} integer infinitely
often, with probability $1$.
\end{theorem}

\begin{whybox}
The standard tool is the \emph{expected number of visits}. Let
$N = \sum_{n=0}^\infty \ind\{S_n = 0\}$ count how many times (including time
$0$) the walk sits at the origin. By linearity (the indicators need not be
independent),
\[
  \E[N] = \sum_{n=0}^\infty \P(S_n = 0).
\]
There is a beautiful dichotomy: the walk is recurrent if and only if
$\E[N] = \infty$. Why? If returns are certain, the walk returns again and
again forever, so $N = \infty$ surely and $\E[N] = \infty$. If there is a
chance $1 - u > 0$ of never returning, then the number of returns is geometric
with finite mean, so $\E[N] < \infty$. So we just have to decide whether
$\sum_n \P(S_n = 0)$ converges.

Return to $0$ is only possible after an even number of steps, with $n/2$ rights
and $n/2$ lefts. Setting $n = 2m$,
\[
  \P(S_{2m} = 0) = \binom{2m}{m}\Big(\tfrac{1}{2}\Big)^{2m}.
\]
By Stirling's formula, $\binom{2m}{m} \sim 4^m/\sqrt{\pi m}$, so
\[
  \P(S_{2m} = 0) \sim \frac{1}{\sqrt{\pi m}}.
\]
Therefore
\[
  \E[N] = \sum_{m=0}^\infty \P(S_{2m}=0)
  \sim \sum_{m=1}^\infty \frac{1}{\sqrt{\pi m}} = \infty,
\]
because $\sum 1/\sqrt{m}$ diverges (it is a $p$-series with $p = 1/2 \le 1$).
The expected number of visits is infinite, so the walk is recurrent. \qed
\end{whybox}

This calculation deserves a long pause, because the \emph{exact same machine}
decides recurrence in every dimension, and the verdict flips on the razor's
edge of how fast $\P(S_{2m}=0)$ decays. In one dimension it decays like
$m^{-1/2}$, the sum diverges, and the walk comes home. Watch what happens as we
add dimensions.

\subsection{Two dimensions: still recurrent}

In two dimensions the walker stands at a lattice point and, each step, moves
north, south, east, or west with probability $1/4$ each. To return to the
origin after $2m$ steps, the net horizontal displacement and the net vertical
displacement must both be zero. A clean way to see the count: think of the
$x$-coordinate and $y$-coordinate as \emph{nearly} two independent 1D walks
(after a $45^\circ$ rotation they become exactly independent), so the return
probability roughly squares:
\[
  \P(S_{2m} = 0) \sim \Big(\frac{1}{\sqrt{\pi m}}\Big)^2 = \frac{1}{\pi m}.
\]
Now the series in question is the harmonic series:
\[
  \E[N] = \sum_m \P(S_{2m}=0) \sim \sum_m \frac{1}{\pi m} = \infty.
\]
It \emph{still} diverges -- but just barely, since $\sum 1/m$ is the famous
divergent harmonic series sitting right at the boundary. So the 2D walk is also
recurrent. The drunk in a city grid \emph{does} get home.

\subsection{Three dimensions: transient}

In three dimensions the walker moves along one of six directions
($\pm x, \pm y, \pm z$) with probability $1/6$ each. Now the return probability
behaves like the cube of the 1D rate:
\[
  \P(S_{2m} = 0) \sim \frac{C}{m^{3/2}}
\]
for a constant $C > 0$. And here the series \emph{converges}:
\[
  \E[N] = \sum_m \P(S_{2m}=0) \sim \sum_m \frac{C}{m^{3/2}} < \infty,
\]
a $p$-series with $p = 3/2 > 1$. Finite expected visits means the 3D walk is
\emph{transient}: with positive probability it never returns to the origin.
The bird flying in three-dimensional space may well wander off forever.

\begin{paradoxbox}
\textbf{P\'olya's theorem: ``A drunk man finds his way home, but a drunk bird
may not.''} This memorable line is due to the mathematician Shizuo Kakutani,
summarizing a $1921$ theorem of George P\'olya. The symmetric lattice walk is
recurrent in dimensions $1$ and $2$ but transient in dimension $3$ and higher.

What makes it a genuine paradox is how \emph{close} the calls are. The entire
fate of the walker hinges on whether a single series converges, and the series
exponents are $1/2$, $1$, $3/2$ in dimensions $1$, $2$, $3$. The boundary
between ``always returns'' and ``may escape'' falls exactly between $d = 2$
(harmonic series, just diverges) and $d = 3$ (just converges). Two dimensions
is recurrent by a hair; three is transient by a hair. Geometry, not luck,
decides -- there is simply \emph{more room to get lost} in three dimensions.

A concrete number: the probability that the 3D walk ever returns to the origin
is about $0.3405$. So a drunk bird has roughly a $34\%$ chance of coming home
and a $66\%$ chance of being lost to the void. The drunk man's chance is
$100\%$.
\end{paradoxbox}

\begin{historybox}
George P\'olya later told the story of how he discovered the theorem. While
walking in the woods near Zurich, he kept running into the same young couple --
again and again, no matter which path he took. He wondered: was it pure chance,
or was something about the structure of paths forcing the encounters? That idle
question on a walk became the recurrence theorem of $1921$, one of the
foundational results of the modern theory of random walks and Markov chains.
The lesson is worth keeping: deep mathematics often begins with a concrete,
almost childish question -- ``why do I keep bumping into these people?'' -- taken
seriously.
\end{historybox}

\begin{happeningbox}
Why does dimension matter so much? Recurrence is a competition between two
forces. The walk \emph{spreads} -- after $n$ steps it occupies a cloud of
diameter $\sim \sqrt{n}$. In $d$ dimensions that cloud has volume $\sim
n^{d/2}$ and contains $\sim n^{d/2}$ lattice sites, so the chance of being on
any one particular site (the origin) is $\sim n^{-d/2}$. Summing over all
times asks whether $\sum_n n^{-d/2}$ converges, and that flips at $d = 2$. The
higher the dimension, the more sites the probability is ``diluted'' across, and
the less likely any specific site -- including home -- gets revisited. Recurrence
is really a statement about whether the spreading cloud is dense enough to keep
hitting the same point.
\end{happeningbox}

\section{The Reflection Principle and the Ballot Problem}
\label{ch26:reflection}

So far we have asked about the walk's position at a fixed time. The richer
questions are about the \emph{path as a whole}: How high does it ever get? When
does it first hit a level? How long does it stay positive? The master tool for
all of these is a single, gorgeous idea.

\begin{intuitionbox}
\textbf{The reflection principle, in one picture}
(Figure~\ref{ch26:fig-reflection}). Suppose a path goes from a
point $A$ above some horizontal line to a point $B$ also above the line, but
\emph{touches} the line somewhere in between. Reflect the portion of the path
\emph{after} the first touch across the line. You get a brand-new path from $A$
to $B'$, the mirror image of $B$. This reflection is a perfect bijection
between (paths from $A$ to $B$ that touch the line) and (all paths from $A$ to
$B'$). Counting one set is the same as counting the other -- and the reflected
set is usually far easier to count.
\end{intuitionbox}

\begin{figure}[H]
\centering
\begin{tikzpicture}[scale=0.62]
\draw[->] (-0.5,0) -- (15,0) node[right]{step $n$};
\draw[->] (0,-5.7) -- (0,4.5) node[above]{$S_n$};
\draw[gray!55, thick, dashed] (0,-1) -- (14,-1) node[right]{\small level $-1$};
\draw[sourceblue, very thick]
 (0,0)--(1,1)--(2,0)--(3,1)--(4,0)--(5,-1)--(6,0)--(7,-1)--(8,0)--(9,1)--(10,2)--(11,1)--(12,2)--(13,3);
\fill[black] (0,0) circle (3pt) node[above left]{$A$};
\fill[black] (13,3) circle (3pt) node[right]{$B$};
\fill[densityred] (5,-1) circle (3.5pt);
\node[densityred, below] at (5,-1) {\small first touch};
\draw[densityred, very thick, dashed]
 (5,-1)--(6,-2)--(7,-1)--(8,-2)--(9,-3)--(10,-4)--(11,-3)--(12,-4)--(13,-5);
\fill[densityred] (13,-5) circle (3pt) node[right]{$B'$ (mirror of $B$)};
\end{tikzpicture}
\caption{The reflection principle. A path from $A$ to $B$ that touches the
dashed level is reflected, after its first touch, into a path from $A$ to the
mirror point $B'$. This pairing is exactly one-to-one, turning a hard count
(``paths that touch'') into an easy one (``all paths to $B'$'').}
\label{ch26:fig-reflection}
\end{figure}

Let us put it to work on the classic first-passage count.

\begin{theorem}[Reflection principle / hitting count]
\label{ch26:thm-reflection}
For $a, b > 0$, the number of $n$-step $\pm 1$ paths from height $a$ to height
$b$ that touch level $0$ at some time equals the total number of $n$-step paths
from $-a$ to $b$.
\end{theorem}

\begin{whybox}
Take any path from $a > 0$ to $b > 0$ that touches level $0$. Let $\tau$ be the
first time it hits $0$. Reflect the entire path \emph{before} time $\tau$ across
the level $0$: the starting height $a$ becomes $-a$, the touch point stays at
$0$ (it is on the mirror), and the path after $\tau$ is left unchanged. The
result is a path from $-a$ to $b$. Conversely, any path from $-a$ to $b$ must
cross level $0$ (it starts below and ends above), and reflecting its piece up
to the first crossing returns a path from $a$ that touches $0$. The map is its
own inverse -- a bijection. So the two counts are equal. \qed
\end{whybox}

This bijection is the engine behind a charming old puzzle.

\begin{example}[The ballot problem]
\label{ch26:ex-ballotworked}
In an election, candidate A receives $\alpha$ votes and candidate B receives
$\beta$ votes, with $\alpha > \beta$. The ballots are counted one at a time in
random order. What is the probability that A is \emph{strictly ahead} of B
throughout the entire count -- from the very first vote to the last?

Model the running lead $S_k = (\text{A votes so far}) - (\text{B votes so
far})$ as a walk: an A-ballot is a $+1$ step, a B-ballot is a $-1$ step. After
all $\alpha + \beta$ ballots the walk sits at $\alpha - \beta > 0$. ``A always
strictly ahead'' means $S_k > 0$ for all $k \ge 1$, i.e.\ the walk never
returns to (or dips below) $0$ after the first step.

The first ballot must be an A (else A is not ahead at step $1$). After that,
starting from height $1$, we need a path to $\alpha - \beta$ that never touches
$0$. By the reflection principle the count of paths that \emph{do} touch $0$
equals the count of paths starting from the mirror height $-1$ -- and a short
bookkeeping with binomial coefficients collapses the whole thing into a
startlingly clean answer:
\[
  \boxed{\,\P(\text{A always strictly ahead}) = \frac{\alpha - \beta}{\alpha + \beta}.\,}
\]
\end{example}

\begin{aestheticbox}
Look at how clean that is. The probability that A leads wire-to-wire depends
only on the \emph{margin} $\alpha - \beta$ relative to the \emph{total}
$\alpha + \beta$ -- the individual values barely matter. If A wins $600$ to
$400$, the chance she led at every single moment of the count is exactly
$200/1000 = 1/5$, no matter how large the electorate. The answer is so simple
that it cries out for a slick proof, and indeed there are several (Andr\'e's
reflection argument above; a cyclic-shift ``cycle lemma'' of Dvoretzky and
Motzkin). When an answer is this clean, nature is telling you that you have
found the right way to count.
\end{aestheticbox}

\begin{historybox}
The ballot problem was posed and solved by Joseph Bertrand in $1887$, and
D\'esir\'e Andr\'e published the elegant reflection argument the same year --
which is why the technique is often called \emph{Andr\'e's reflection
method}. It was one of the first deep results about the \emph{trajectory} of a
random process rather than its endpoint, and it foreshadowed the modern theory
of fluctuations: the arcsine laws, the distribution of the maximum, and
first-passage times all descend from this single reflection idea.
\end{historybox}

\section{Gambler's Ruin}
\label{ch26:ruin}

The reflection principle counts paths that avoid a barrier. The gambler's ruin
problem asks the complementary, intensely practical question: if there are
barriers on \emph{both} sides -- you stop when you go broke or when you reach a
target -- which barrier do you hit, and how long does it take?

\begin{walkbox}
A gambler starts with $\$i$ and bets $\$1$ per round on a game that pays even
money. She quits if she reaches a fortune of $\$N$ (her goal) or if she hits
$\$0$ (ruin). With a fair game, what is the chance she reaches her goal before
going broke? And how many rounds, on average, until the game ends?
\end{walkbox}

This is a random walk on the integers $0, 1, 2, \dots, N$ with two
\emph{absorbing} barriers: once you hit $0$ or $N$ the walk stops. Let $p_i$ be
the probability of reaching $N$ before $0$, starting from $i$.

\begin{theorem}[Gambler's ruin]
\label{ch26:thm-ruin}
For the fair game ($p = q = 1/2$), starting from $i$,
\[
  \P(\text{reach } N \text{ before } 0) = \frac{i}{N}.
\]
For a biased game with $p \ne q$ and ratio $r = q/p$,
\[
  \P(\text{reach } N \text{ before } 0) = \frac{1 - r^{\,i}}{1 - r^{\,N}}.
\]
The expected duration of the fair game, starting from $i$, is
\[
  \E[\text{number of rounds}] = i\,(N - i).
\]
\end{theorem}

\begin{whybox}
Condition on the first step -- this is the ``first-step analysis'' that solves
almost every absorbing-walk problem. From state $i$ (with $0 < i < N$) the walk
moves to $i+1$ with probability $p$ or to $i-1$ with probability $q$, so
\[
  p_i = p\,p_{i+1} + q\,p_{i-1}, \qquad p_0 = 0,\ p_N = 1.
\]
For the fair case $p = q = 1/2$ this says $p_i = \tfrac{1}{2}(p_{i+1} +
p_{i-1})$, i.e.\ each $p_i$ is the \emph{average} of its neighbors. A sequence
whose every term is the average of its neighbors is a straight line. With
boundary values $p_0 = 0$ and $p_N = 1$, the only such line is $p_i = i/N$.

For the biased case, rearrange into $p(p_{i+1} - p_i) = q(p_i - p_{i-1})$; the
consecutive differences $d_i = p_i - p_{i-1}$ satisfy $d_{i+1} = r\,d_i$ with
$r = q/p$, so $d_i = d_1 r^{i-1}$. Summing the geometric series of differences
and applying $p_0 = 0$, $p_N = 1$ gives $p_i = (1 - r^i)/(1 - r^N)$.

For the expected duration $D_i$, a parallel first-step analysis gives
$D_i = 1 + \tfrac{1}{2}D_{i+1} + \tfrac{1}{2}D_{i-1}$ with $D_0 = D_N = 0$
(the $+1$ counts the current round). The function $D_i = i(N-i)$ satisfies both
the recurrence and the boundaries: $i(N-i)$ is a downward parabola, and one
checks $1 + \tfrac12(i{+}1)(N{-}i{-}1) + \tfrac12(i{-}1)(N{-}i{+}1) = i(N-i)$.
\qed
\end{whybox}

\begin{example}[The house always wins -- slowly]
\label{ch26:ex-ruin}
You sit down at a fair-coin game with $\$50$, hoping to double to $\$100$
(so $i = 50$, $N = 100$). Your chance of success is $p_i = 50/100 = 1/2$ --
perfectly fair, as it must be. And the expected number of rounds is $i(N-i) =
50 \times 50 = 2500$. A long evening.

Now suppose the game is slightly unfair, as every real casino game is. Take
$p = 0.49$, $q = 0.51$, so $r = q/p = 0.51/0.49 \approx 1.0408$. Then
\[
  p_{50} = \frac{1 - 1.0408^{50}}{1 - 1.0408^{100}}
  \approx \frac{1 - 7.39}{1 - 54.63} \approx 0.12.
\]
A house edge of a mere $2\%$ collapses your chance of doubling up from $50\%$ to
about $12\%$. This is the brutal arithmetic of the casino: a tiny per-bet bias
compounds, through $r^N$, into near-certain ruin for anyone who plays long
enough.
\end{example}

Figure~\ref{ch26:fig-ruin} plots the success probability against the starting
stake for both a fair and a slightly unfair game, and the gap between the two
curves is the whole story of the casino.

\begin{figure}[H]
\centering
\begin{tikzpicture}
\begin{axis}[
  width=12cm, height=6.2cm,
  xlabel={starting stake $i$ (goal $N=100$)},
  ylabel={$\P(\text{reach } N \text{ before } 0)$},
  xmin=0, xmax=100, ymin=0, ymax=1,
  grid=both, grid style={gray!18},
  legend style={font=\scriptsize, at={(0.02,0.98)}, anchor=north west},
  title={Gambler's ruin: chance of reaching the goal vs.\ starting stake},
]
\addplot[sourceblue, very thick, domain=0:100, samples=2] {x/100};
\addplot[densityred, very thick] coordinates {
(0,0.0000)(5,0.0041)(10,0.0092)(15,0.0153)(20,0.0229)(25,0.0320)(30,0.0433)(35,0.0570)(40,0.0737)(45,0.0942)(50,0.1192)(55,0.1497)(60,0.1870)(65,0.2325)(70,0.2881)(75,0.3560)(80,0.4390)(85,0.5404)(90,0.6641)(95,0.8153)(100,1.0000)
};
\legend{fair game ($p=\tfrac12$): $i/N$, biased ($p=0.49$): $(1-r^i)/(1-r^N)$}
\end{axis}
\end{tikzpicture}
\caption{Success probability in gambler's ruin. For a fair game it is the
straight line $i/N$. A mere $2\%$ house edge ($p=0.49$, $r=q/p\approx1.041$)
bends the curve sharply downward: starting at the halfway stake $i=50$ your
chance of doubling falls from $50\%$ to about $12\%$, matching
Example~\ref{ch26:ex-ruin}.}
\label{ch26:fig-ruin}
\end{figure}

\begin{paradoxbox}
\textbf{The bold-play paradox.} Given an unfair game and a fixed goal, what
betting strategy maximizes your chance of reaching it? Intuition says ``bet
small, survive longer, give yourself more chances.'' Intuition is exactly
wrong. The longer you expose your money to a negative-edge game, the more
surely the edge grinds you down -- duration is the enemy. The optimal strategy
is \emph{bold play}: bet as much as you can each time, racing to the goal in as
few rounds as possible to minimize the house's chances to apply its edge.
Against an unfair game, the rational move is to gamble \emph{harder}, not
softer. (Against a \emph{fair} game it makes no difference -- timid and bold play
both leave your success probability at $i/N$.)
\end{paradoxbox}

\begin{connectionbox}
Gambler's ruin is your first encounter with a finite \emph{Markov chain} with
absorbing states, and the ``each value is the average of its neighbors''
equation $p_i = \tfrac12(p_{i+1}+p_{i-1})$ is the discrete \emph{Laplace
equation}. Its solutions are \emph{harmonic functions}; the absorbing barriers
play the role of boundary conditions. This is the same equation that governs
steady-state heat and electrostatic potential, and it is exactly why random
walks turn out to be a probabilistic method for solving partial differential
equations. We develop the Markov-chain machinery -- transition matrices,
absorption probabilities, expected hitting times -- in the next chapter; the
linear-algebra view of these systems lives in the companion \emph{Linear
Algebra From the Ground Up}.
\end{connectionbox}

\section{From Random Walk to Brownian Motion}
\label{ch26:brownian}

We end where the random walk has been pointing all along: at the continuous
limit. What happens if we take smaller and smaller steps, faster and faster, in
just the right balance?

\begin{intuitionbox}
Zoom out from a random-walk path until the individual $\pm 1$ steps blur into a
continuous wiggle. To keep the picture from collapsing to a flat line or
exploding to infinity, we must rescale space and time together -- and the
$\sqrt{n}$ law tells us exactly how. If we take $n$ steps in one unit of time,
the walk spreads by $\sqrt{n}$; to keep that spread finite as $n \to \infty$ we
must shrink each step by a factor of $\sqrt{n}$. Space scales like the square
root of time. The limiting object is \emph{Brownian motion}.
\end{intuitionbox}

Make it precise. Rescale the walk by letting $n$ steps fit into time $t$ and
shrinking each step to size $1/\sqrt{n}$:
\[
  B_t^{(n)} = \frac{1}{\sqrt{n}}\, S_{\lfloor nt \rfloor}.
\]
At time $t = 1$ this is $S_n/\sqrt{n}$, which by the CLT converges to
$\Normal(0, 1)$. More is true: the entire rescaled \emph{path} converges (this
is Donsker's theorem, the ``functional CLT'') to a random continuous function
$B_t$ with these properties.

\begin{definition}[Brownian motion]
\label{ch26:def-brownian}
\emph{Standard Brownian motion} $\{B_t : t \ge 0\}$ is the random process with:
$B_0 = 0$; independent increments (disjoint time intervals give independent
displacements); stationary Gaussian increments, $B_t - B_s \sim \Normal(0,
t-s)$ for $s < t$; and continuous sample paths.
\end{definition}

Notice the fingerprints of the random walk all over this definition.
Independent increments are the i.i.d.\ steps. The variance $t - s$ growing
linearly in elapsed time is the $\Var(S_n) = n$ law, survived into the limit.
The Gaussian shape is the CLT doing its work on each increment. Brownian
motion is the simple random walk, viewed from infinitely far away.
Figure~\ref{ch26:fig-brownian} shows one such limiting path.

\begin{figure}[H]
\centering
\begin{tikzpicture}
\begin{axis}[
  width=13cm, height=6cm,
  xlabel={time $t$}, ylabel={$B_t$},
  xmin=0, xmax=1, ymin=-1.7, ymax=1.7,
  grid=both, grid style={gray!18},
  title={A Brownian sample path as the scaling limit of a random walk},
]
\addplot[sourceblue, thick] coordinates {
(0,0)(0.02,0.14)(0.04,0.05)(0.06,0.22)(0.08,0.10)(0.10,0.31)(0.12,0.20)(0.14,0.41)(0.16,0.33)(0.18,0.52)(0.20,0.40)(0.22,0.55)(0.24,0.47)(0.26,0.66)(0.28,0.58)(0.30,0.45)(0.32,0.60)(0.34,0.74)(0.36,0.63)(0.38,0.80)(0.40,0.70)(0.42,0.58)(0.44,0.72)(0.46,0.61)(0.48,0.50)(0.50,0.64)(0.52,0.55)(0.54,0.42)(0.56,0.30)(0.58,0.44)(0.60,0.33)(0.62,0.20)(0.64,0.08)(0.66,-0.05)(0.68,0.10)(0.70,-0.02)(0.72,-0.16)(0.74,-0.04)(0.76,-0.20)(0.78,-0.33)(0.80,-0.22)(0.82,-0.38)(0.84,-0.50)(0.86,-0.40)(0.88,-0.55)(0.90,-0.44)(0.92,-0.60)(0.94,-0.48)(0.96,-0.63)(0.98,-0.52)(1.0,-0.66)
};
\addplot[black, dashed, samples=60, domain=0:1]{sqrt(x)};
\addplot[black, dashed, samples=60, domain=0:1]{-sqrt(x)};
\end{axis}
\end{tikzpicture}
\caption{A sample path of Brownian motion on $[0,1]$. The dashed curves
$\pm\sqrt{t}$ mark one standard deviation; the path is continuous everywhere
yet -- as a deeper analysis shows -- differentiable nowhere, infinitely jagged at
every scale, exactly as the un-zoomable wiggle of the underlying random walk
demands.}
\label{ch26:fig-brownian}
\end{figure}

\begin{historybox}
In $1827$ the botanist Robert Brown peered through a microscope at pollen grains
suspended in water and saw them jiggling ceaselessly, with no apparent cause.
For nearly eighty years the ``Brownian movement'' resisted explanation. Then in
his \emph{annus mirabilis} of $1905$, Albert Einstein argued that the jiggle was
the pollen being bombarded by countless invisible water molecules -- a visible
random walk driven by molecular collisions -- and derived that the mean squared
displacement grows \emph{linearly} in time, the continuous echo of $\Var(S_n) =
n$. This gave Jean Perrin the means to measure Avogadro's number and clinched
the physical reality of atoms. The rigorous mathematical object was constructed
by Norbert Wiener in $1923$, which is why Brownian motion is also called the
\emph{Wiener process}. A drunkard's stagger and a pollen grain's tremble obey
the very same law.
\end{historybox}

\begin{applicationbox}
\textbf{Brownian motion in the wild.} The random-walk-to-Brownian-motion limit
is one of the most useful pipelines in applied mathematics.
\begin{itemize}
\item \textbf{Finance.} The Black--Scholes model treats the logarithm of a
stock price as Brownian motion with drift. Option pricing, hedging, and modern
quantitative finance all rest on this foundation -- the daily up/down ticks of a
price are the random-walk steps.
\item \textbf{Physics and chemistry.} Diffusion of heat, ink, perfume, and
pollutants is Brownian motion of enormous numbers of particles; the
$\sqrt{t}$ spreading law sets how fast a drop of dye fills a glass of water.
\item \textbf{Genetics.} The frequency of a neutral allele in a finite
population performs a random walk (the Wright--Fisher model), whose
diffusion limit predicts how long until a gene is fixed or lost.
\item \textbf{Machine learning.} Stochastic gradient descent is a noisy walk
through a loss landscape; its continuous-time idealization is a Brownian-driven
diffusion, and that lens explains why small-batch noise can help models
generalize.
\item \textbf{Algorithms.} Randomized search, Markov chain Monte Carlo (MCMC),
and many randomized algorithms are analyzed as walks on enormous state spaces.
\end{itemize}
\end{applicationbox}

\begin{labbox}
Simulate it yourself. Generate $n = 1000$ i.i.d.\ $\pm 1$ steps and plot the
cumulative sum -- you will see a jagged path like
Figure~\ref{ch26:fig-paths}. Now do it for $100$ independent walks and, at each
time $n$, compute the standard deviation across the $100$ endpoints; plot that
against $\sqrt{n}$ and watch them track. Finally rescale one long walk by
dividing positions by $\sqrt{n}$ and squeezing time into $[0,1]$: the result is
a convincing Brownian path. Three short loops turn this whole chapter into a
picture on your screen.
\end{labbox}

\begin{recoverybox}
\textbf{The thread of the chapter.} A random walk is a running sum of i.i.d.\
steps. Its mean stays put ($\E[S_n] = 0$) but its spread grows like $\sqrt{n}$
($\Var(S_n) = n$). Whether it returns home is decided by a single series and
flips with dimension -- recurrent in 1D and 2D, transient in 3D. The reflection
principle turns hard path-counts into easy ones, giving the ballot problem and
first-passage results for free. With absorbing barriers we get gambler's ruin,
solved by first-step analysis and the discrete Laplace equation. And scaling
the walk down in space and time, guided by the $\sqrt{n}$ law, produces
Brownian motion -- the same process that jiggles a pollen grain and prices an
option. Every one of these is a question about that one wandering token we
started with.
\end{recoverybox}

\section{Exercises}
\label{ch26:exercises}

\begin{exercise}[Warm-up positions]
\label{ch26:ex-warmup}
For the simple symmetric random walk, compute $\P(S_4 = 2)$ and $\P(S_4 = 0)$
directly by listing or counting the $\pm 1$ step sequences. Confirm your
answers against the binomial formula $\P(S_n = k) = \binom{n}{(n+k)/2}2^{-n}$.
\end{exercise}

\begin{exercise}[Mean and variance, biased]
\label{ch26:ex-biasedmoments}
A biased walk has $\P(X_i = +1) = 0.6$. Find $\E[S_{100}]$ and
$\SD(S_{100})$. Using the normal approximation, estimate $\P(S_{100} > 40)$.
\end{exercise}

\begin{exercise}[The $\sqrt{n}$ law in numbers]
\label{ch26:ex-sqrtlaw}
For the symmetric walk, how many steps must you take before the standard
deviation of your position first exceeds $50$? Before it exceeds $500$? Explain
in one sentence why the second answer is one hundred times the first, not ten.
\end{exercise}

\begin{exercise}[Return probability]
\label{ch26:ex-return}
Compute $\P(S_{2m} = 0)$ exactly for $m = 1, 2, 3$ for the symmetric walk, and
compare each with the Stirling approximation $1/\sqrt{\pi m}$. Comment on how
fast the approximation kicks in.
\end{exercise}

\begin{exercise}[Recurrence vs.\ transience, conceptual]
\label{ch26:ex-rectrans}
Explain, in your own words and without heavy formulas, why
$\sum_n \P(S_n = 0) = \infty$ implies the walk returns to $0$ with probability
$1$, while a finite sum implies positive escape probability. Then state which
$p$-series exponent governs each dimension $d = 1, 2, 3$ and why $d = 2$ is the
borderline case.
\end{exercise}

\begin{exercise}[Ballot problem]
\label{ch26:ex-ballot}
Candidate A gets $7$ votes and B gets $3$. (a) Using the ballot theorem, find
the probability A is strictly ahead throughout the count. (b) Find the
probability A is \emph{never behind} (ties allowed). \emph{Hint for (b): the
``never behind'' probability is given by the closely related formula
$(\alpha - \beta + 1)/(\alpha + 1)$.}
\end{exercise}

\begin{exercise}[Gambler's ruin, fair]
\label{ch26:ex-ruinfair}
You start with $\$30$ and play a fair $\$1$ game, quitting at $\$0$ or
$\$100$. (a) What is your probability of reaching $\$100$? (b) What is the
expected number of rounds? (c) If instead you start at $\$50$, by what factor
does the expected duration change, and why is it not simply proportional to
your starting capital?
\end{exercise}

\begin{exercise}[Gambler's ruin, biased]
\label{ch26:ex-ruinbiased}
With $p = 0.45$, $q = 0.55$ (so $r = q/p = 11/9$), you start with $\$10$ and
aim for $\$20$. Compute your probability of success using
$p_i = (1 - r^i)/(1 - r^N)$. Compare with the fair-game value $i/N = 1/2$ and
comment on how a $10\%$ edge per bet translates into the overall outcome.
\end{exercise}

\begin{exercise}[First-step analysis from scratch]
\label{ch26:ex-firststep}
Consider a walk on states $0, 1, 2, 3$ with absorbing barriers at $0$ and $3$,
fair steps in between. Write down the equations $p_i = \tfrac12 p_{i+1} +
\tfrac12 p_{i-1}$ with $p_0 = 0$, $p_3 = 1$, and solve the small linear system
by hand for $p_1$ and $p_2$. Verify they match $i/N$.
\end{exercise}

\begin{exercise}[Scaling to Brownian motion]
\label{ch26:ex-scaling}
Define $B_t^{(n)} = S_{\lfloor nt \rfloor}/\sqrt{n}$. (a) Show that
$\E[B_t^{(n)}] = 0$ and $\Var(B_t^{(n)}) \to t$ as $n \to \infty$. (b) For
$s < t$, argue that $B_t^{(n)} - B_s^{(n)}$ is asymptotically independent of
$B_s^{(n)}$ and has variance approaching $t - s$. Identify which property of
Brownian motion each part establishes.
\end{exercise}

\begin{exercise}[Proof: maximum of the walk]
\label{ch26:ex-maxproof}
Use the reflection principle to prove that for the symmetric walk and any
integer $b > 0$,
\[
  \P\!\Big(\max_{0 \le k \le n} S_k \ge b\Big)
  = 2\,\P(S_n > b) + \P(S_n = b).
\]
\emph{Hint: a path reaching level $b$ either ends at or above $b$, or ends
below $b$ after touching it; reflect the latter to pair it with a path ending
above $b$.}
\end{exercise}

\begin{exercise}[Expected duration is a parabola]
\label{ch26:ex-durparabola}
Verify directly that $D_i = i(N - i)$ solves the duration recurrence
$D_i = 1 + \tfrac12 D_{i+1} + \tfrac12 D_{i-1}$ with $D_0 = D_N = 0$. Then
explain intuitively why the game lasts longest when you start in the middle
($i = N/2$) and is shortest near a barrier.
\end{exercise}
